\documentclass[10pt,a4paper]{amsart}
\usepackage[margin=19mm]{geometry}
\usepackage{amsmath,amssymb,mathtools}
\usepackage[scr=rsfs,cal=euler]{mathalfa}
\usepackage{xfrac}
\usepackage[unicode,hidelinks,bookmarks=false]{hyperref}
\newcommand{\ie}{i.e.\ }
\newcommand{\cf}{cf.\ }
\newcommand{\resp}{respectively\ }
\newcommand{\Subsection}{\subsection}
\providecommand{\nobreakdash}{\nobreak\hbox{-}}
\setlength{\emergencystretch}{2em}
\multlinegap0pt
\usepackage{amssymb}
\usepackage{xcolor}
\newtheorem{theorem}{Theorem}
\title{A Rigorous Proof of a Ramanujan Machine Identity for $-\pi/4$ via Exact Recurrence Solving}
\author{Chao Wang}
\date{Source: arXiv:2601.08461v2 (CC BY 4.0)}
\begin{document}
\maketitle

\noindent\emph{Source and licence.} A Rigorous Proof of a Ramanujan Machine Identity for $-\pi/4$ via Exact Recurrence Solving. Version arXiv:2601.08461v2 (CC BY 4.0), distributed by arXiv under the Creative Commons Attribution 4.0 International licence, http://creativecommons.org/licenses/by/4.0/, as recorded in the arXiv licence metadata for this exact version and shown on the abstract page https://arxiv.org/abs/2601.08461v2. Full licence notice: \texttt{licenses/arxiv-2601.08461-CC-BY-4.0.txt}.\par\smallskip

\noindent\textit{Editorial scope.} Counted unit: the article's theorem thm:closed\_form (source0.tex lines 135--143). The article's complete original proof of it is delivered with it (source0.tex lines 145--177, one \texttt{proof} environment of 33 lines). No mathematics is written, repaired or re-derived: the statement and the proof are the article's own lines, in the article's own notation, and no retained line is substituted or re-flowed.  Retained uncounted as the source-local context the proof needs: lines 84--88; lines 91--97.  Nothing was omitted from any retained span, so the package carries no omission record.  No retained line contains a citation, so the wrapper carries no bibliography and no bibliography entry.  No retained line contains a figure environment, an image-inclusion command or drawing code, so no image is delivered and the package declares no asset dependency.  Editorial formatting only, in addition: an AMS \texttt{amsart} preamble that restates the article's own declarations and macros used by the retained lines, namely the environments \texttt{theorem} on separate counters, together with the standard packages those lines need.  The retained environments print in this document as Theorem 1 in order of appearance, whereas the article prints the same items with the numbering of its own full document.  The complete native source of the article as distributed in the arXiv e-print for this exact version is bundled unchanged as \texttt{sources/192597/source0.tex}.
  \begin{equation}\label{eq:conjecture}
    -\frac{\pi}{4}
    = \cfrac{1}{-1 + \cfrac{1}{-4 + \cfrac{-2}{-7 + \cfrac{-9}{-10
          + \cfrac{-20}{-13 + \cdots}}}}},
  \end{equation}

  \begin{equation}\label{eq:an_def}
    a_n =
    \begin{cases}
      1,               & n = 1, 2,     \\
      -(n-1)(2n-5),    & n \ge 3.
    \end{cases}
  \end{equation}

\begin{theorem}[Closed Form of Denominators]\label{thm:closed_form}
  Let the denominator sequence $q_n$ be defined by the recurrence
  $q_n = b_n q_{n-1} + a_n q_{n-2}$ with initial values $q_{-1}=0,\,q_0=1$.
  For the coefficients given in~\eqref{eq:conjecture} and~\eqref{eq:an_def},
  the exact solution for all $n\ge 0$ is
  \begin{equation}\label{eq:qn_closed}
    q_n = (-1)^n\,(2n-3)!!\,(n^2+n-1).
  \end{equation}
\end{theorem}

\begin{proof}
  We proceed by strong induction on $n$.

  \paragraph{Base cases.}
  \begin{itemize}
  \item $n=0$: The formula gives $(-1)^0(-3)!!(-1) = 1\cdot(-1)\cdot(-1)
    = 1 = q_0$.
  \item $n=1$: The formula gives $(-1)^1(-1)!!(1) = -1$. The recurrence
    gives $q_1 = b_1 q_0 = (-1)(1) = -1$. Matches.
  \item $n=2$: The formula gives $(-1)^2(1)!!(5) = 5$. The recurrence
    gives $q_2 = b_2 q_1 + a_2 q_0 = (-4)(-1)+(1)(1) = 5$. Matches.
  \end{itemize}

  \paragraph{Inductive step.}
  Assume~\eqref{eq:qn_closed} holds for all indices less than $n$, where
  $n\ge 3$. With $b_n = -(3n-2)$ and $a_n = -(n-1)(2n-5)$, we substitute
  the inductive hypothesis into the recurrence. Setting
  $T_n = (-1)^n(2n-5)!!$ and using $(-1)^{n-1} = -(-1)^n$,
  $(-1)^{n-2} = (-1)^n$, and $(2n-5)(2n-7)!! = (2n-5)!!$:
  \begin{align*}
    q_n &= -(3n-2)\bigl[(-1)^{n-1}(2n-5)!!\,(n^2-n-1)\bigr]
           -(n-1)(2n-5)\bigl[(-1)^{n-2}(2n-7)!!\,(n^2-3n+1)\bigr] \\
        &= T_n\bigl[(3n-2)(n^2-n-1) - (n-1)(n^2-3n+1)\bigr].
  \end{align*}
  Expanding: $(3n-2)(n^2-n-1) = 3n^3-5n^2-n+2$ and
  $(n-1)(n^2-3n+1) = n^3-4n^2+4n-1$, so the bracket equals
  $2n^3-n^2-5n+3$. The target formula with
  $(2n-3)!! = (2n-3)(2n-5)!!$ gives
  \begin{equation*}
    q_n = T_n(2n-3)(n^2+n-1) = T_n(2n^3-n^2-5n+3),
  \end{equation*}
  which coincides with the recurrence evaluation. The induction is complete.
\end{proof}

\end{document}
